\documentclass[12pt,a4paper]{article}

\usepackage[utf8]{inputenc}
\usepackage[vietnamese]{babel}
\usepackage{amsmath,amssymb,amsthm}
\usepackage{geometry}
\usepackage{enumitem}
\usepackage[most]{tcolorbox}
\usepackage{hyperref}
\usepackage{fancyhdr}
\usepackage{tikz}
\usepackage{eso-pic}
\usepackage{xcolor}

\geometry{a4paper, margin=2cm, bottom=2.5cm}
\hypersetup{colorlinks=true, linkcolor=blue!70!black, urlcolor=blue}
\setlist[itemize]{topsep=4pt, itemsep=2pt}
\setlist[enumerate]{topsep=4pt, itemsep=2pt}

\AddToShipoutPictureFG{%
  \AtPageCenter{%
    \begin{tikzpicture}[remember picture, overlay]
      \node[rotate=45, scale=1, opacity=0.06]{\fontsize{60}{72}\selectfont\bfseries\sffamily Đặng Tấn Quí};
    \end{tikzpicture}%
  }%
}

\pagestyle{fancy}
\fancyhf{}
\fancyhead[L]{\small\textit{Các chủ đề nâng cao}}
\fancyhead[R]{\small\textit{Đặng Tấn Quí}}
\fancyfoot[C]{\thepage}
\fancyfoot[L]{\footnotesize\textcopyright\ Đặng Tấn Quí}
\fancyfoot[R]{\footnotesize SĐT: 0377\,122\,210}
\renewcommand{\headrulewidth}{0.4pt}
\renewcommand{\footrulewidth}{0.4pt}
\fancypagestyle{plain}{\fancyhf{}\fancyfoot[C]{\thepage}\fancyfoot[L]{\footnotesize\textcopyright\ Đặng Tấn Quí}\fancyfoot[R]{\footnotesize SĐT: 0377\,122\,210}\renewcommand{\headrulewidth}{0pt}\renewcommand{\footrulewidth}{0.4pt}}

\newtcolorbox{dinhnghia}[1][]{enhanced, breakable, colback=blue!3!white, colframe=blue!60!black, fonttitle=\bfseries, title={Định nghĩa}, #1}
\newtcolorbox{dinhly}[1][]{enhanced, breakable, colback=green!3!white, colframe=green!50!black, fonttitle=\bfseries, title={Định lý}, #1}
\newtcolorbox{tinhchat}[1][]{enhanced, breakable, colback=yellow!5!white, colframe=orange!50!black, fonttitle=\bfseries, title={Tính chất}, #1}
\newtcolorbox{phuongphap}[1][]{enhanced, breakable, colback=gray!5!white, colframe=gray!50!black, fonttitle=\bfseries, title={Phương pháp}, #1}
\newtcolorbox{luuy}[1][]{enhanced, breakable, colback=red!3!white, colframe=red!50!black, fonttitle=\bfseries, title={Lưu ý}, #1}
\newtcolorbox{congthuc}[1][]{enhanced, breakable, colback=cyan!3!white, colframe=cyan!50!black, fonttitle=\bfseries, title={Công thức}, #1}
\newtcolorbox{vidu}[1][]{enhanced, breakable, colback=violet!3!white, colframe=violet!50!black, fonttitle=\bfseries, title={Ví dụ}, #1}

\begin{document}

\thispagestyle{empty}
\begin{tikzpicture}[remember picture, overlay]
  \fill[blue!80!black] (current page.south west) rectangle (current page.north east);
  \fill[blue!60!cyan, opacity=0.8] (current page.south west) -- (current page.north west) -- ([xshift=-3cm]current page.north east) -- (current page.south east) -- cycle;
  \foreach \r/\o in {6/0.05, 4.5/0.08, 3/0.12} {\fill[white, opacity=\o] ([xshift=2cm, yshift=-2cm]current page.north east) circle (\r cm);}
  \draw[white, line width=2pt] ([yshift=-5cm]current page.north west) ++(2cm,0) -- ++(17cm,0);
  \draw[white, line width=2pt] ([yshift=-16cm]current page.north west) ++(2cm,0) -- ++(17cm,0);
  \node[anchor=west, white, font=\fontsize{26}{32}\bfseries\selectfont] at ([xshift=3cm, yshift=-7.5cm]current page.north west) {PHẦN 6: CÁC CHỦ ĐỀ};
  \node[anchor=west, white, font=\fontsize{26}{32}\bfseries\selectfont] at ([xshift=3cm, yshift=-9.2cm]current page.north west) {NÂNG CAO};
  \node[anchor=west, white, font=\fontsize{16}{20}\selectfont] at ([xshift=3cm, yshift=-11cm]current page.north west) {Tài liệu luyện thi du học};
  \node[white, opacity=0.10, font=\fontsize{80}{80}\selectfont, rotate=-10] at ([xshift=-3cm, yshift=4cm]current page.center) {$\sum$};
  \node[anchor=west, white, font=\Large\bfseries] at ([xshift=3cm, yshift=-13cm]current page.north west) {Biên soạn:};
  \node[anchor=west, yellow!80!white, font=\fontsize{28}{34}\bfseries\selectfont] at ([xshift=3cm, yshift=-14.5cm]current page.north west) {ĐẶNG TẤN QUÍ};
  \node[anchor=west, white!90, font=\large] at ([xshift=3cm, yshift=-18cm]current page.north west) {SĐT: 0377\,122\,210};
  \node[anchor=south, white!80, font=\small] at ([yshift=1.5cm]current page.south) {Tài liệu có bản quyền --- Cấm sao chép khi chưa được phép};
  \node[anchor=south east, white, font=\Large\bfseries] at ([xshift=-2cm, yshift=3cm]current page.south east) {\the\year};
\end{tikzpicture}
\newpage

\thispagestyle{empty}
\vspace*{5cm}
\begin{center}
  {\Large\bfseries PHẦN 6: CÁC CHỦ ĐỀ NÂNG CAO}\\[8pt]
  {\large Tài liệu luyện thi du học}
\end{center}
\vspace{2cm}
\begin{tcolorbox}[enhanced, colback=red!3!white, colframe=red!60!black, fonttitle=\bfseries, title={Thông tin bản quyền}, boxrule=1.5pt]
  \textbf{Tác giả:} Đặng Tấn Quí \quad \textbf{Liên hệ:} 0377\,122\,210 \quad \textbf{Năm:} \the\year \\[6pt]
  \textbf{Bản quyền \textcopyright\ \the\year\ Đặng Tấn Quí. Bảo lưu mọi quyền.}
\end{tcolorbox}
\vfill\begin{center}\small\textit{Tài liệu lưu hành nội bộ}\end{center}
\newpage
\tableofcontents
\newpage

%% ================================================================
%%  PHẦN 6: CÁC CHỦ ĐỀ NÂNG CAO
%% ================================================================
\section{Các chủ đề nâng cao}

%% ----------------------------------------------------------------
%%  6.1 Phương trình tham số
%% ----------------------------------------------------------------
\subsection{Phương trình tham số (Parametric Equations): Đạo hàm, độ dài cung}

\begin{dinhnghia}[title={Đường cong tham số}]
  Đường cong được mô tả bởi \textbf{phương trình tham số} là:
  \[
    x = f(t), \quad y = g(t), \quad t \in [a, b].
  \]
  Tham số $t$ thường là thời gian. Khi $t$ tăng từ $a$ đến $b$, điểm $(x(t), y(t))$ vẽ ra đường cong.
\end{dinhnghia}

\begin{congthuc}[title={Đạo hàm theo trục $x$}]
  \[
    \frac{dy}{dx} = \frac{dy/dt}{dx/dt} = \frac{g'(t)}{f'(t)} \quad (f'(t) \ne 0).
  \]
  Đạo hàm bậc hai:
  \[
    \frac{d^2y}{dx^2} = \frac{d}{dx}\!\left[\frac{dy}{dx}\right] = \frac{\dfrac{d}{dt}\!\left[\dfrac{dy}{dx}\right]}{dx/dt}.
  \]
\end{congthuc}

\begin{congthuc}[title={Độ dài cung (Arc Length)}]
  \[
    L = \int_a^b \sqrt{\left(\frac{dx}{dt}\right)^2 + \left(\frac{dy}{dt}\right)^2}\,dt.
  \]
\end{congthuc}

\begin{congthuc}[title={Diện tích mặt tròn xoay (Surface Area)}]
  Quay đường cong tham số quanh trục $Ox$:
  \[
    S = 2\pi \int_a^b |y(t)| \sqrt{\left(\frac{dx}{dt}\right)^2 + \left(\frac{dy}{dt}\right)^2}\,dt.
  \]
\end{congthuc}

\begin{vidu}
  Cho đường cycloid $x = t - \sin t$, $y = 1 - \cos t$, $t \in [0, 2\pi]$.
  \begin{enumerate}[label=(\alph*)]
    \item Tìm $\dfrac{dy}{dx}$ tại $t = \dfrac{\pi}{2}$.
    \item Tính độ dài cung một cung cycloid.
  \end{enumerate}

  \medskip\textbf{Giải:}
  \begin{enumerate}[label=(\alph*)]
    \item $\dfrac{dx}{dt} = 1 - \cos t$, $\dfrac{dy}{dt} = \sin t$.
      $\dfrac{dy}{dx} = \dfrac{\sin t}{1 - \cos t}$.
      Tại $t = \dfrac{\pi}{2}$: $\dfrac{dy}{dx} = \dfrac{1}{1-0} = 1$.

    \item $L = \displaystyle\int_0^{2\pi} \sqrt{(1-\cos t)^2 + \sin^2 t}\,dt = \int_0^{2\pi} \sqrt{2 - 2\cos t}\,dt$.
      Dùng $2-2\cos t = 4\sin^2(t/2)$:
      $L = \displaystyle\int_0^{2\pi} 2|\sin(t/2)|\,dt = 2\int_0^{2\pi} \sin(t/2)\,dt = 2\left[-2\cos(t/2)\right]_0^{2\pi} = 8$.
  \end{enumerate}
\end{vidu}

%% ----------------------------------------------------------------
%%  6.2 Tọa độ cực
%% ----------------------------------------------------------------
\subsection{Tọa độ cực (Polar Coordinates): Đạo hàm và tính diện tích miền cực}

\begin{dinhnghia}[title={Tọa độ cực}]
  Mỗi điểm trong mặt phẳng được biểu diễn bởi $(r, \theta)$:
  \begin{itemize}
    \item $r$: khoảng cách từ gốc tọa độ (cực) đến điểm.
    \item $\theta$: góc tạo bởi tia kéo đến điểm và trục cực (dương $= $ trục $Ox$).
  \end{itemize}

  \textbf{Chuyển đổi:}
  \[
    x = r\cos\theta, \quad y = r\sin\theta, \quad r^2 = x^2 + y^2, \quad \tan\theta = \frac{y}{x}.
  \]
\end{dinhnghia}

\begin{congthuc}[title={Đạo hàm trong tọa độ cực}]
  Với $r = f(\theta)$ (xem như đường cong tham số với $x = r\cos\theta$, $y = r\sin\theta$):
  \[
    \frac{dy}{dx} = \frac{f'(\theta)\sin\theta + f(\theta)\cos\theta}{f'(\theta)\cos\theta - f(\theta)\sin\theta}.
  \]
\end{congthuc}

\begin{congthuc}[title={Diện tích miền cực}]
  Diện tích vùng giới hạn bởi $r = f(\theta)$ với $\alpha \le \theta \le \beta$:
  \[
    S = \frac{1}{2}\int_\alpha^\beta [f(\theta)]^2\,d\theta = \frac{1}{2}\int_\alpha^\beta r^2\,d\theta.
  \]
  Diện tích giữa hai đường cong cực $r_1(\theta) \le r_2(\theta)$:
  \[
    S = \frac{1}{2}\int_\alpha^\beta \bigl[r_2(\theta)^2 - r_1(\theta)^2\bigr]\,d\theta.
  \]
\end{congthuc}

\begin{congthuc}[title={Độ dài cung trong tọa độ cực}]
  \[
    L = \int_\alpha^\beta \sqrt{r^2 + \left(\frac{dr}{d\theta}\right)^2}\,d\theta.
  \]
\end{congthuc}

\begin{vidu}
  Tính diện tích vùng giới hạn bởi cardioid $r = 1 + \cos\theta$.

  \medskip\textbf{Giải:}
  \[
    S = \frac{1}{2}\int_0^{2\pi} (1+\cos\theta)^2\,d\theta
    = \frac{1}{2}\int_0^{2\pi} (1 + 2\cos\theta + \cos^2\theta)\,d\theta.
  \]
  Dùng $\cos^2\theta = \dfrac{1+\cos 2\theta}{2}$:
  \[
    = \frac{1}{2}\int_0^{2\pi} \left(\frac{3}{2} + 2\cos\theta + \frac{\cos 2\theta}{2}\right)\,d\theta
    = \frac{1}{2}\left[\frac{3\theta}{2} + 2\sin\theta + \frac{\sin 2\theta}{4}\right]_0^{2\pi}
    = \frac{1}{2} \cdot 3\pi = \frac{3\pi}{2}.
  \]
\end{vidu}

%% ----------------------------------------------------------------
%%  6.3 Vectơ trong mặt phẳng
%% ----------------------------------------------------------------
\subsection{Vectơ trong mặt phẳng: Vị trí, vận tốc, gia tốc}

\begin{dinhnghia}[title={Vectơ vị trí, vận tốc, gia tốc}]
  Cho vật chuyển động với vectơ vị trí $\vec{r}(t) = \langle x(t),\, y(t) \rangle$ (hoặc trong không gian $\langle x, y, z \rangle$):
  \begin{itemize}
    \item \textbf{Vận tốc:} $\vec{v}(t) = \vec{r}'(t) = \langle x'(t),\, y'(t) \rangle$.
    \item \textbf{Tốc độ (speed):} $|\vec{v}(t)| = \sqrt{[x'(t)]^2 + [y'(t)]^2}$.
    \item \textbf{Gia tốc:} $\vec{a}(t) = \vec{v}'(t) = \vec{r}''(t) = \langle x''(t),\, y''(t) \rangle$.
  \end{itemize}
\end{dinhnghia}

\begin{congthuc}[title={Quãng đường và độ dời}]
  \begin{itemize}
    \item \textbf{Độ dời (displacement)} từ $t=a$ đến $t=b$: $\vec{r}(b) - \vec{r}(a) = \displaystyle\int_a^b \vec{v}(t)\,dt$.
    \item \textbf{Quãng đường (distance traveled)}: $s = \displaystyle\int_a^b |\vec{v}(t)|\,dt$.
    \item Tọa độ tại thời điểm $t$: $\vec{r}(t) = \vec{r}(0) + \displaystyle\int_0^t \vec{v}(u)\,du$.
  \end{itemize}
\end{congthuc}

\begin{vidu}
  Vật có vận tốc $\vec{v}(t) = \langle \cos t,\, \sin t \rangle$, ban đầu tại $(0,0)$.
  \begin{enumerate}[label=(\alph*)]
    \item Tìm vị trí tại thời điểm $t$.
    \item Tính quãng đường từ $t=0$ đến $t = 2\pi$.
  \end{enumerate}

  \medskip\textbf{Giải:}
  \begin{enumerate}[label=(\alph*)]
    \item $\vec{r}(t) = \langle \sin t, -\cos t + 1 \rangle$ (vật chạy trên đường tròn $x^2 + (y-1)^2 = 1$).
    \item $|\vec{v}(t)| = \sqrt{\cos^2 t + \sin^2 t} = 1$. Quãng đường $= \displaystyle\int_0^{2\pi} 1\,dt = 2\pi$.
  \end{enumerate}
\end{vidu}

%% ----------------------------------------------------------------
%%  6.4 Dãy số và tiêu chuẩn hội tụ
%% ----------------------------------------------------------------
\subsection{Dãy số vô hạn và các tiêu chuẩn hội tụ}

\begin{dinhnghia}[title={Dãy số và giới hạn}]
  Dãy số $\{a_n\}$ hội tụ đến $L$ nếu $\displaystyle\lim_{n\to\infty} a_n = L$ (tồn tại và hữu hạn). Ngược lại là phân kỳ.
\end{dinhnghia}

\begin{dinhnghia}[title={Chuỗi số (Series)}]
  Chuỗi số là tổng vô hạn $\displaystyle\sum_{n=1}^\infty a_n$. Tổng riêng thứ $N$ là $S_N = \displaystyle\sum_{n=1}^N a_n$.

  Chuỗi \textbf{hội tụ} nếu $\displaystyle\lim_{N\to\infty} S_N = S$ (tồn tại hữu hạn); ngược lại \textbf{phân kỳ}.
\end{dinhnghia}

\begin{tinhchat}[title={Các chuỗi đặc biệt}]
  \begin{itemize}
    \item \textbf{Chuỗi hình học:} $\displaystyle\sum_{n=0}^\infty ar^n = \dfrac{a}{1-r}$ \; hội tụ khi $|r| < 1$; phân kỳ khi $|r| \ge 1$.
    \item \textbf{Chuỗi $p$:} $\displaystyle\sum_{n=1}^\infty \dfrac{1}{n^p}$: hội tụ khi $p > 1$; phân kỳ khi $p \le 1$.
    \item \textbf{Chuỗi harmonic:} $\displaystyle\sum_{n=1}^\infty \dfrac{1}{n}$ phân kỳ (trường hợp $p=1$).
  \end{itemize}
\end{tinhchat}

\begin{congthuc}[title={Các tiêu chuẩn hội tụ}]
  \textbf{1. Tiêu chuẩn so sánh (Comparison Test):}
  Nếu $0 \le a_n \le b_n$: $\sum b_n$ hội tụ $\Rightarrow$ $\sum a_n$ hội tụ; $\sum a_n$ phân kỳ $\Rightarrow$ $\sum b_n$ phân kỳ.

  \textbf{2. Tiêu chuẩn so sánh giới hạn (Limit Comparison):}
  Nếu $\displaystyle\lim_{n\to\infty} \dfrac{a_n}{b_n} = c > 0$: $\sum a_n$ và $\sum b_n$ cùng hội tụ hoặc cùng phân kỳ.

  \textbf{3. Tiêu chuẩn tích phân (Integral Test):}
  Nếu $f$ giảm, dương, liên tục trên $[1,\infty)$ và $f(n) = a_n$: $\sum a_n$ hội tụ $\Leftrightarrow$ $\displaystyle\int_1^\infty f(x)\,dx$ hội tụ.

  \textbf{4. Tiêu chuẩn tỉ số (Ratio Test):}
  $L = \displaystyle\lim_{n\to\infty}\left|\dfrac{a_{n+1}}{a_n}\right|$: hội tụ nếu $L < 1$; phân kỳ nếu $L > 1$; không kết luận nếu $L = 1$.

  \textbf{5. Tiêu chuẩn căn (Root Test):}
  $L = \displaystyle\lim_{n\to\infty} \sqrt[n]{|a_n|}$: hội tụ nếu $L < 1$; phân kỳ nếu $L > 1$; không kết luận nếu $L = 1$.

  \textbf{6. Tiêu chuẩn Leibniz (Alternating Series Test):}
  Chuỗi xen dấu $\displaystyle\sum (-1)^n b_n$ ($b_n > 0$) hội tụ nếu $b_n$ giảm và $b_n \to 0$.
\end{congthuc}

\begin{vidu}
  Xét sự hội tụ của $\displaystyle\sum_{n=1}^\infty \dfrac{n^2}{n!}$ bằng tiêu chuẩn tỉ số.

  \medskip\textbf{Giải:}
  $\dfrac{a_{n+1}}{a_n} = \dfrac{(n+1)^2}{(n+1)!} \cdot \dfrac{n!}{n^2} = \dfrac{(n+1)^2}{(n+1) \cdot n^2} = \dfrac{n+1}{n^2} \to 0$ khi $n \to \infty$.

  Vì $L = 0 < 1$, chuỗi \textbf{hội tụ tuyệt đối}.
\end{vidu}

%% ----------------------------------------------------------------
%%  6.5 Chuỗi lũy thừa
%% ----------------------------------------------------------------
\subsection{Chuỗi lũy thừa (Power Series) và bán kính hội tụ}

\begin{dinhnghia}[title={Chuỗi lũy thừa}]
  Chuỗi lũy thừa tâm $a$ là:
  \[
    \sum_{n=0}^\infty c_n (x-a)^n = c_0 + c_1(x-a) + c_2(x-a)^2 + \cdots
  \]
\end{dinhnghia}

\begin{dinhly}[title={Bán kính và khoảng hội tụ}]
  Với mỗi chuỗi lũy thừa, tồn tại $R \in [0, +\infty]$ gọi là \textbf{bán kính hội tụ} sao cho:
  \begin{itemize}
    \item Chuỗi hội tụ tuyệt đối khi $|x - a| < R$.
    \item Chuỗi phân kỳ khi $|x - a| > R$.
    \item Tại $|x - a| = R$: cần kiểm tra riêng.
  \end{itemize}
  Tính $R$ bằng: $\dfrac{1}{R} = \displaystyle\lim_{n\to\infty}\left|\dfrac{c_{n+1}}{c_n}\right|$ hoặc $\dfrac{1}{R} = \lim_{n\to\infty} \sqrt[n]{|c_n|}$.
\end{dinhly}

\begin{congthuc}[title={Đạo hàm và tích phân chuỗi lũy thừa}]
  Trong khoảng hội tụ $|x-a| < R$:
  \begin{itemize}
    \item \textbf{Đạo hàm:} $\displaystyle\left[\sum_{n=0}^\infty c_n(x-a)^n\right]' = \sum_{n=1}^\infty n\,c_n(x-a)^{n-1}$.
    \item \textbf{Tích phân:} $\displaystyle\int \sum_{n=0}^\infty c_n(x-a)^n\,dx = \sum_{n=0}^\infty \frac{c_n}{n+1}(x-a)^{n+1} + C$.
  \end{itemize}
  Bán kính hội tụ không thay đổi sau khi đạo hàm hoặc tích phân.
\end{congthuc}

\begin{vidu}
  Tìm bán kính và khoảng hội tụ của $\displaystyle\sum_{n=1}^\infty \dfrac{(-1)^n x^n}{n}$.

  \medskip\textbf{Giải:}
  $\dfrac{1}{R} = \displaystyle\lim_{n\to\infty}\left|\dfrac{c_{n+1}}{c_n}\right| = \lim_{n\to\infty} \dfrac{n}{n+1} = 1 \Rightarrow R = 1$.

  Kiểm tra đầu mút:
  \begin{itemize}
    \item $x = 1$: $\displaystyle\sum \dfrac{(-1)^n}{n}$ -- chuỗi xen dấu, hội tụ (Leibniz).
    \item $x = -1$: $\displaystyle\sum \dfrac{-1}{n}$ -- chuỗi harmonic âm, phân kỳ.
  \end{itemize}
  Khoảng hội tụ: $(-1, 1]$.
\end{vidu}

%% ----------------------------------------------------------------
%%  6.6 Chuỗi Taylor và Maclaurin
%% ----------------------------------------------------------------
\subsection{Chuỗi Taylor và chuỗi Maclaurin (Xấp xỉ hàm số bằng đa thức)}

\begin{dinhnghia}[title={Chuỗi Taylor và Maclaurin}]
  Nếu $f$ có đạo hàm mọi bậc tại $a$, chuỗi Taylor của $f$ tâm $a$ là:
  \[
    f(x) = \sum_{n=0}^\infty \frac{f^{(n)}(a)}{n!}(x-a)^n = f(a) + f'(a)(x-a) + \frac{f''(a)}{2!}(x-a)^2 + \cdots
  \]
  \textbf{Chuỗi Maclaurin} là chuỗi Taylor tâm $a = 0$:
  \[
    f(x) = \sum_{n=0}^\infty \frac{f^{(n)}(0)}{n!} x^n.
  \]
\end{dinhnghia}

\begin{congthuc}[title={Các chuỗi Maclaurin quan trọng}]
  \begin{align*}
    e^x &= \sum_{n=0}^\infty \frac{x^n}{n!} = 1 + x + \frac{x^2}{2!} + \frac{x^3}{3!} + \cdots & (\text{mọi } x) \\[4pt]
    \sin x &= \sum_{n=0}^\infty \frac{(-1)^n x^{2n+1}}{(2n+1)!} = x - \frac{x^3}{3!} + \frac{x^5}{5!} - \cdots & (\text{mọi } x) \\[4pt]
    \cos x &= \sum_{n=0}^\infty \frac{(-1)^n x^{2n}}{(2n)!} = 1 - \frac{x^2}{2!} + \frac{x^4}{4!} - \cdots & (\text{mọi } x) \\[4pt]
    \frac{1}{1-x} &= \sum_{n=0}^\infty x^n = 1 + x + x^2 + \cdots & (|x| < 1) \\[4pt]
    \ln(1+x) &= \sum_{n=1}^\infty \frac{(-1)^{n+1} x^n}{n} = x - \frac{x^2}{2} + \frac{x^3}{3} - \cdots & (-1 < x \le 1) \\[4pt]
    (1+x)^k &= \sum_{n=0}^\infty \binom{k}{n} x^n & (|x|<1)
  \end{align*}
\end{congthuc}

\begin{vidu}
  Tìm chuỗi Maclaurin của $f(x) = e^{-x^2}$ và dùng nó để xấp xỉ $\displaystyle\int_0^1 e^{-x^2}\,dx$ đến 3 chữ số thập phân.

  \medskip\textbf{Giải:}
  Thay $x \to -x^2$ vào chuỗi của $e^x$:
  \[
    e^{-x^2} = \sum_{n=0}^\infty \frac{(-1)^n x^{2n}}{n!} = 1 - x^2 + \frac{x^4}{2!} - \frac{x^6}{3!} + \cdots
  \]
  Tích phân hạng theo hạng từ $0$ đến $1$:
  \[
    \int_0^1 e^{-x^2}\,dx = 1 - \frac{1}{3} + \frac{1}{10} - \frac{1}{42} + \frac{1}{216} - \cdots \approx 0.747.
  \]
\end{vidu}

%% ----------------------------------------------------------------
%%  6.7 Sai số Lagrange
%% ----------------------------------------------------------------
\subsection{Sai số Lagrange (Lagrange Error Bound)}

\begin{dinhly}[title={Phần dư Taylor và cận sai số Lagrange}]
  Khi dùng đa thức Taylor bậc $n$ ($T_n(x)$) để xấp xỉ $f(x)$ tại điểm $x$ (tâm $a$):
  \[
    R_n(x) = f(x) - T_n(x) = \frac{f^{(n+1)}(c)}{(n+1)!}(x-a)^{n+1}
  \]
  cho một số $c$ nào đó giữa $a$ và $x$ (định lý Lagrange).

  \textbf{Cận trên của sai số:}
  \[
    |R_n(x)| \le \frac{M}{(n+1)!}|x-a|^{n+1},
  \]
  trong đó $M = \max_{t \text{ giữa } a \text{ và } x} |f^{(n+1)}(t)|$.
\end{dinhly}

\begin{phuongphap}[title={Sử dụng Lagrange Error Bound}]
  \begin{enumerate}
    \item Xác định $n$ (bậc đa thức đang dùng).
    \item Tính $f^{(n+1)}(x)$ và tìm $M = \max |f^{(n+1)}(t)|$ trên đoạn liên quan.
    \item Thế vào công thức: $|R_n(x)| \le \dfrac{M}{(n+1)!}|x-a|^{n+1}$.
  \end{enumerate}
\end{phuongphap}

\begin{vidu}
  Dùng đa thức Taylor bậc 3 của $\sin x$ tâm $0$ để xấp xỉ $\sin(0.1)$. Ước lượng sai số.

  \medskip\textbf{Giải:}
  $T_3(x) = x - \dfrac{x^3}{6}$. Xấp xỉ: $\sin(0.1) \approx 0.1 - \dfrac{0.001}{6} \approx 0.09983$.

  Phần dư bậc 4: $f^{(4)}(x) = \sin x$, $|f^{(4)}(x)| \le 1$ với mọi $x$.

  $|R_3(0.1)| \le \dfrac{1}{4!}|0.1|^4 = \dfrac{0.0001}{24} \approx 4.17 \times 10^{-6}$.

  Rất chính xác (sai số không quá $5 \times 10^{-6}$).
\end{vidu}

\medskip\noindent\textbf{Các dạng bài tập tổng hợp:}
\begin{enumerate}[label=\textbf{Dạng \arabic*:}, leftmargin=*, align=left]
  \item Tính đạo hàm và độ dài cung của đường cong tham số.
  \item Tính diện tích trong tọa độ cực.
  \item Xác định sự hội tụ của dãy và chuỗi số.
  \item Tìm bán kính và khoảng hội tụ của chuỗi lũy thừa.
  \item Tìm chuỗi Taylor/Maclaurin của hàm số.
  \item Ước lượng sai số bằng Lagrange Error Bound.
\end{enumerate}

\end{document}
